\documentclass[11pt]{article}

\usepackage[margin=1in]{geometry}
\usepackage{booktabs}
\usepackage{hyperref}
\usepackage{enumitem}
\usepackage{xcolor}
\usepackage{array}

\setlength{\emergencystretch}{3em}

\title{Harnessing the Vibe: Benchmarking Intervention-Driven Harness Recovery for Coding Agents}
\author{Anonymous Author(s)}
\date{}

\newcommand{\todo}[1]{\textcolor{red}{[TODO: #1]}}

\begin{document}
\maketitle

\begin{abstract}
Coding agents increasingly operate as semi-autonomous software engineers, but
interactive ``vibe coding'' workflows expose a recurring reliability gap: when an
agent fails, the user often fixes the immediate symptom manually, while the
underlying harness limitation remains unmodeled and the final result is not
replayed from a clean state. We argue that such episodes are not merely task
failures; they are intervention-triggered harness failures. This paper proposes
VibeHarness, a supervisor-worker protocol that checkpoints the repository,
delegates implementation to worker agents, observes intervention evidence,
writes a falsifiable decision contract, patches the harness, rolls back, and
replays the task until both the application change and the harness repair are
verified. We also introduce VibeHarnessBench, a benchmark design for evaluating
whether coding-agent systems can convert human interventions into durable,
replayable harness improvements. The benchmark scores task success,
intervention absorption, attribution accuracy, regression-free replay,
trajectory safety, episode completeness, and transfer of harness fixes across
related tasks.
\todo{Replace this abstract with empirical results after running the benchmark.}
\end{abstract}

\section{Introduction}

Repository-level coding agents are commonly evaluated by whether their final
patch passes tests. This view misses a failure mode that is common in real
interactive usage. A user asks for a feature, the agent gets stuck, the user
manually starts a service, opens a browser, clarifies acceptance criteria, or
points out a missed regression, and the agent eventually completes the task. The
session may look successful, but the system has learned nothing durable: the next
task can fail in the same way, and the final state may never be verified from a
clean checkpoint.

We study the hypothesis that many apparent coding-agent failures are recoverable
harness failures. A harness is the executable environment around the model:
prompts, tools, setup scripts, memories, review procedures, validation commands,
and coordination protocols. Recent work has recognized harnesses as central to
agent performance, but interactive user interventions remain under-specified as
an evaluation signal. We focus on the moment when a human intervention reveals a
missing harness capability.

The intended contribution of this work is threefold:

\begin{enumerate}[leftmargin=*]
  \item a taxonomy of intervention-triggered harness failures in vibe coding;
  \item VibeHarnessBench, a replay-based benchmark for harness recovery;
  \item VibeHarness, a supervisor-worker protocol that repairs the harness,
  rolls back, and validates the task from the original checkpoint;
  \item an episode-package scoring model that audits not only final patch
  success but also attribution, permissions, information flow, and transfer.
\end{enumerate}

\section{Motivating Example}

Consider a web application task: ``Add an empty-state message to the todo list.''
The worker edits the component and passes unit tests. The user opens the app and
reports that the empty-state message overlaps the first todo item after clicking
Add. The immediate problem is a UI bug, but the deeper failure is that the
harness gave the agent no browser validation affordance. A robust workflow should
not simply patch the component after the user report. It should add or repair the
browser-testing harness, roll back to the checkpoint before the task, and replay
the original request to verify that the same user intervention is no longer
needed.

\section{Related Work}

\paragraph{Software-engineering agent benchmarks.}
SWE-bench introduced repository-level issue resolution from real GitHub
issues~\cite{jimenez2024swebench}. SWE-bench Verified attempted to curate a more
reliable subset, but later analysis reported test-quality and contamination
problems severe enough that OpenAI recommended moving beyond it for frontier
coding evaluation~\cite{openai2026swebenchverified}. PatchDiff further showed
that patches counted as solved can still diverge behaviorally from developer
expectations~\cite{wang2025solved}. These findings motivate evaluation beyond a
single final pass/fail signal.

\paragraph{Vibe-coding benchmarks.}
FeatBench evaluates feature-level coding from natural language requirements and
finds that agents often make aggressive changes that cause regressions
\cite{chen2025featbench}. SUSVIBES studies security vulnerabilities in
agent-generated feature implementations~\cite{zhao2025susvibes}.
SWE-WebDevBench evaluates AI app builders as virtual software agencies and
reports gaps in specification handling, engineering quality, and production
readiness~\cite{saxena2026swewebdevbench}. Vibe Code Bench evaluates
end-to-end web application development with browser-based workflows and
human-aligned evaluators~\cite{tran2026vibecodebench}. These benchmarks
strengthen the case for realistic interaction-level evaluation, but they do not
directly score whether user interventions become durable harness improvements.

\paragraph{Harness engineering.}
OpenAI's harness-engineering report describes workflows where Codex agents use
standard development tools, browser automation, observability, review agents, and
repository-readable artifacts to complete product work~\cite{openai2026harness}.
Agentic Harness Engineering proposes observability-driven automatic evolution of
coding-agent harnesses and reports gains on Terminal-Bench and SWE-bench
Verified~\cite{lin2026ahe}. Code as Agent Harness surveys the broader shift
toward executable, verifiable, stateful agent systems~\cite{ning2026codeharness}.
AgentFlow synthesizes multi-agent harnesses for vulnerability discovery using a
typed graph DSL and diagnostic feedback~\cite{liu2026agentflow}. Our focus is
narrower: the supervised recovery loop triggered by user intervention in
interactive coding tasks.

AI Harness Engineering formalizes a runtime substrate with responsibilities such
as task specification, context selection, tool access, memory, task state,
observability, failure attribution, verification, permissions, entropy auditing,
and intervention recording~\cite{zhong2026aiharness}. HarnessAudit shows why
output-level evaluation is insufficient: a harness may produce a correct final
answer while violating permission or information-flow constraints during the
trajectory~\cite{liu2026harnessaudit}. VibeHarness adopts these insights by
requiring episode packages and trajectory audits.

Natural-Language Agent Harnesses further argues that reusable harness policy can
be made executable and inspectable as natural-language artifacts
\cite{pan2026nlah}. This supports VibeHarness's repository-readable contracts
and adapter files.

\paragraph{Agent-system maintenance.}
AgentIssue-Bench studies whether agents can fix issues in agent systems and
shows low resolution rates, indicating that agent infrastructure itself is a hard
maintenance target~\cite{rahardja2025agentissues}. VibeHarnessBench differs by
scoring whether an intervention trace is converted into a replayable harness
repair, rather than only whether an agent-system bug is fixed.

\paragraph{Recovery and optimization.}
Wink studies production coding-agent misbehaviors and lightweight
self-intervention, reporting reductions in engineer interventions and tool-call
failures~\cite{nanda2026wink}. VeRO introduces versioning, rewards,
observations, and structured traces for agents that optimize other agents
\cite{ursekar2026vero}. VibeHarness is closest to this line of work, but it
focuses on converting explicit user interventions into durable harness artifacts
and requiring replay from the original checkpoint.

\section{VibeHarness Protocol}

VibeHarness separates task implementation from harness recovery.

\begin{enumerate}[leftmargin=*]
  \item \textbf{Checkpoint.} Before task execution, commit or otherwise snapshot
  the repository, dependencies, and harness state.
  \item \textbf{Delegate.} A supervisor gives the original user request to one or
  more worker agents.
  \item \textbf{Observe.} The supervisor records failures, tool traces, user
  interventions, validation gaps, permission events, and state transitions.
  \item \textbf{Attribute.} The supervisor classifies the failure as a task bug
  or harness gap and localizes it to runtime components such as tool access,
  task specification, observability, verification, or coordination.
  \item \textbf{Contract.} Before editing the harness, the supervisor writes a
  decision contract: the hypothesized cause, intended harness edit, expected
  effect, and falsification check.
  \item \textbf{Repair harness.} If the problem is a harness gap, the supervisor
  changes only allowed harness artifacts such as bootstrap scripts, tool
  instructions, tests, memories, or review protocols.
  \item \textbf{Rollback.} The repository returns to the pre-task checkpoint,
  retaining only the repaired harness where permitted.
  \item \textbf{Replay and verify.} The worker reruns the original task. The run
  succeeds only if hidden task checks and harness replay checks pass without new
  human intervention.
  \item \textbf{Audit.} The run produces an episode package containing the
  checkpoint id, harness version, role graph, intervention trace, attribution,
  decision contract, harness diff, replay proof, and permission audit.
\end{enumerate}

\section{Design Principles}

Table~\ref{tab:adopted-principles} summarizes how VibeHarness absorbs the most
solid parts of prior work while preserving a narrower contribution.

\begin{table}[ht]
\centering
\small
\begin{tabular}{>{\raggedright\arraybackslash}p{0.25\linewidth}>{\raggedright\arraybackslash}p{0.31\linewidth}>{\raggedright\arraybackslash}p{0.34\linewidth}}
\toprule
Prior work & Solid idea adopted & VibeHarness adaptation \\
\midrule
Agentic Harness Engineering & Component, experience, and decision
observability; edits as falsifiable contracts & Every harness edit names the
component, evidence, prediction, and replay falsifier. \\
AgentFlow & Typed multi-agent harness graph over roles, tools, channels, and
coordination & The supervisor-worker-critic-auditor topology is represented as
a typed graph with declared permissions. \\
Code as Agent Harness & Code artifacts as executable, inspectable, stateful
substrate & Bootstrap, browser, acceptance, memory, and review artifacts become
versioned harness code. \\
OpenAI Codex harness report & Make apps, logs, browser behavior, and reviews
readable to agents & Episode packages include app observations, logs,
screenshots, review comments, and replay proof. \\
AI Harness Engineering & Runtime responsibilities for specification, tools,
state, verification, permissions, entropy, and intervention recording & The
benchmark schema requires implicated harness components and audit artifacts. \\
HarnessAudit & Final success can hide trajectory-level violations & VH-score
requires permission and information-flow compliance. \\
\bottomrule
\end{tabular}
\caption{Prior-work ideas absorbed into the VibeHarness framework.}
\label{tab:adopted-principles}
\end{table}

\subsection{Runtime Components}

VibeHarness models a harness as a set of editable runtime components: task
specification, context selection, tool access, project memory, task state,
observability, failure attribution, verification, permissions, entropy auditing,
intervention recording, and coordination. A benchmark task must name which
components are implicated. This prevents a vague claim such as ``the harness was
bad'' and supports component-level transfer analysis.

\subsection{Episode Packages}

The episode package is the unit of evidence. It contains the initial checkpoint,
the harness version, the role graph, tool calls, tool outputs, tests, browser or
log observations, human interventions, attribution, decision contract, harness
diff, rollback command, replay result, and trajectory-safety audit. A run that
passes hidden tests but cannot produce an episode package is scored as
incomplete.

\section{VibeHarnessBench}

Each benchmark task contains a repository fixture, a natural-language user
request, a transfer group, a known harness gap, implicated harness components,
observability requirements, an intervention trace, a decision contract, safety
constraints, allowed harness edit paths, visible checks, hidden checks, a replay
protocol, and required episode artifacts. The benchmark targets six failure
classes: environment bootstrap, tool affordance, verification blind spot, state
drift, specification capture, and coordination breakdown.

\subsection{Metrics}

The primary metric is VH-score: the fraction of tasks where the final patch
passes hidden checks, the harness patch is valid, replay from the original
checkpoint succeeds, no fresh human intervention is required, the trajectory
respects permissions and information-flow boundaries, and the episode package is
complete. Diagnostic metrics include task pass rate, harness replay pass rate,
intervention absorption rate, attribution accuracy, component-localization
accuracy, regression-free rate, human interrupt count, rollback efficiency,
decision-contract pass rate, trajectory-safety pass rate, episode completeness,
and transfer gain.

\subsection{Baselines}

We propose comparing VibeHarness against single-agent execution, single-agent
retry without harness edits, multi-agent review without harness repair, and a
supervisor that edits the harness but does not replay from the checkpoint. Two
additional ablations isolate the framework's stricter claims: a supervisor that
repairs and replays without decision contracts, and a supervisor that repairs and
replays without trajectory auditing.
\todo{Add implementation details and model versions after experiments.}

\section{Experimental Plan}

The initial benchmark should include both synthetic fixtures and real traces.
Synthetic fixtures provide deterministic scoring for early iteration. Real traces
should be collected from coding-agent sessions where a user had to manually
intervene. Each trace should be de-identified, labeled by failure class, and
converted into an executable fixture with visible and hidden checks.

The key empirical questions are:

\begin{enumerate}[leftmargin=*]
  \item How often are user interventions attributable to harness gaps rather
  than application-code difficulty?
  \item Does checkpointed harness recovery improve final task pass rate compared
  with retry and review baselines?
  \item Do harness repairs transfer to held-out tasks in the same failure class?
  \item What is the cost in tokens, wall-clock time, and rollbacks?
\end{enumerate}

\section{Threats to Validity}

Harness recovery can overfit to the benchmark if hidden checks are too similar to
visible intervention evidence. Real traces may leak private project details or
reflect user behavior that is not representative. Replay success depends on
deterministic environments, which is difficult for browser and service-heavy
tasks. The benchmark may also reward excessive harness growth unless entropy
audits penalize scope creep. Finally, the boundary between task code and harness
code can be ambiguous; the benchmark must enforce allowed edit scopes and hidden
information boundaries.

\section{Conclusion}

Vibe coding turns users into informal reliability engineers. When their
interventions are treated as one-off fixes, agent workflows repeatedly fail in
the same way. VibeHarness reframes those moments as data for harness recovery:
observe the intervention, repair the harness, roll back, and replay. The next
step is empirical: build VibeHarnessBench fixtures, run baselines, and measure
whether intervention-driven harness repair improves reliability beyond ordinary
task retries.

\section*{Acknowledgements}

This project builds on recent work in coding-agent benchmarks, harness
engineering, recovery systems, trajectory audit, and production agent tooling. We
thank the authors and maintainers of these systems for making the problem
concrete enough to engineer against. VibeHarness should be read as a narrow
engineering synthesis around intervention-driven harness recovery, not as a
claim to have invented harness engineering.

\bibliographystyle{plain}
\bibliography{references}

\end{document}
